\begin{table}[t]
\centering
\small
\setlength{\tabcolsep}{6pt}
\begin{tabular}{lrr}
\hline
Metric & Baseline & Oracle Rewrite \\
\hline
Entity-ambiguous text retrieval steps & 964 & 132 \\
Wasted-step count & 633 & 71 \\
Wasted-step ratio & 65.66\% & 53.79\% \\
$\Delta F = 0$ count & 915 & 121 \\
$\Delta F = 0$ ratio & 94.92\% & 91.67\% \\
\hline
\end{tabular}
\caption{Stagnation-oriented statistics on text retrieval steps explicitly labeled as entity-ambiguous. Wasted-step ratio is defined as the fraction of entity-ambiguous text queries whose query similarity exceeds the similarity threshold. $\Delta F = 0$ ratio is the fraction of entity-ambiguous text retrieval steps that contribute no fact gain. Oracle rewrite reduces both the wasted-step ratio and the zero-gain ratio, but both remain high, supporting the interpretation that entity ambiguity primarily manifests as trajectory stagnation rather than productive error-correcting exploration.}
\label{tab:entity-ambiguity-stagnation}
\end{table}

\begin{table}[t]
\centering
\small
\setlength{\tabcolsep}{6pt}
\begin{tabular}{lrr}
\hline
$\Delta F = 0$ reason & Baseline & Oracle Rewrite \\
\hline
similarity\_gate & 633 & 71 \\
no\_contribution & 267 & 42 \\
other & 15 & 8 \\
\hline
\end{tabular}
\caption{Reason breakdown among entity-ambiguous text retrieval steps with $\Delta F = 0$. In both settings, the dominant zero-gain source is the similarity gate, which is consistent with repeated or near-duplicate queries and therefore with trajectory stagnation.}
\label{tab:entity-ambiguity-stagnation-breakdown}
\end{table}
